\section{Conclusion}
\label{sec:conclusion}

HMRS connects the state used to manage collections, documents and segments
to the interfaces that agents use to find and inspect evidence. The
collection catalogue and document shelf expose metadata for narrowing
search, and opened documents and indexed segments carry the provenance
needed to inspect a result. Lifecycle checks make discrepancies observable:
a rename left duplicate vector entries, and re-indexing restored the expected
state.

On the 72-question evaluation, the tiered configuration that combined
document and chunk windows recovered more annotated targets than the
baseline tiered configuration. The difference was 15 questions in the
frozen run and 12 in each rerun on the same questions.
This improvement came with a higher text budget than the baseline and
remained well below the budget of reading the corpus in full. The verifier
replay exposed a complementary trade-off. Fusion order placed a gold
document in the top five for 15 of the 16 candidate pools, compared with
seven under verifier-score ordering.
These results motivate two concrete practices in managed document systems:
audit the state that supplies retrieval candidates, and test evidence
admission and ordering as separate decisions.
